\chapterpage{shock}{用冲击波检测工具跟踪波前}
\subhead{从位置已知的密度边缘开始}
从练习包中打开 \code{moving-edge/fort.10}。60 个网格中心等距排列，间隔为 1 µm。密度从约 1 g/cm³ 平滑升至 1.5 g/cm³，过渡区每个输出时刻向右移动一格。文件共有 13 个时刻，覆盖 0 至 12 ns。第 6 ns 的边缘中心在 24.5 µm。
\begin{enumerate}
\item 在\ui{热图}页选择质量密度 \code{R}、欧拉网格中心 \code{XC}、长度 \code{um}、时间 \code{ns} 和线性色标，点击自动范围和绘图。启动检测时，工具会读取主热图当前的变量、坐标和单位。
\item 打开\ui{工具 → 冲击波检测}（\keys{Ctrl+D}）。在\ui{参考时刻}中选第 6 步，搜索范围分别填 \code{20} 和 \code{29}，点击\ui{添加参考}。确认上方表格出现第 6 步和该区间。范围使用热图的位置单位。
\item 将\ui{最大位置偏差}设为 5 µm，\ui{密度变化阈值}保留 0.2，勾选\ui{显示轨迹}，点击\ui{开始检测}。将可视化时刻切到第 6 步，查看这一帧。
\end{enumerate}
\Figure{shock-map}{Ubuntu 中文界面中的检测结果。热图上的轨迹沿着密度过渡区移动，左侧参考表记录了本次跟踪所用的时刻和搜索范围。}{90mm}
\checknote{先核对三个位置}{0、6、12 ns 时，保存的波前位置应依次为 18.5、24.5、30.5 µm。亮色高密度平台还会继续向右延伸。定位边缘时应看密度变化最陡的位置。}

\nextpage{shockinspect}{逐帧核对波前位置与代表密度}
\subhead{参考时刻和检测阈值怎样起作用}
程序先将参考区间内的场值按该区间的最小值与最大值归一化，再找空间梯度绝对值最大的采样位置。区间要覆盖至少三个有效点，且场值须有变化。“密度变化阈值”作用于归一化后的梯度。此后程序向相邻时刻跟踪，用最大位置偏差限制搜索范围。参考表为空时，工具会在所选参考时刻的有效范围内搜索。若同一帧有两个陡边，可在另一时刻添加参考区间，帮助同一条轨迹跟住选定的边缘。
\subhead{把密度曲线和导数曲线对起来看}
选择第 6 步作为可视化时刻，先在\ui{密度分布}页找到上升最陡的区段，再到\ui{导数曲线}页检查峰是否落在 24.5 µm 附近。再看第 0 步和第 12 步。切换可视化时刻只改变正在检查的帧，表中的参考时刻仍为第 6 步。
\Figure{shock-derivative}{第 6 ns 的梯度峰标出边缘中心。选用跳变后密度时，其取样标记在峰的右侧，此处导数的绝对值已降至指定阈值以下。}{83mm}
\subhead{选择与轨迹一同导出的密度}
最大值模式取当前整条剖面上的密度最大值。\ui{Post-shock jump} 模式从整条剖面最强的绝对梯度出发，朝坐标增大的方向找第一个导数绝对值低于阈值的点。把绝对导数阈值设为 \code{0.01} (g/cm³)/µm，第 6 ns 导出的代表密度为 1.496654 g/cm³。最大值模式给出 1.500000 g/cm³。切换模式后重新点击“开始检测”可刷新逐帧检查图。
\checknote{方向与单位}{跳变后取样总是向坐标增大的一侧寻找。数据有多个陡边、密度随位置下降或长度单位改变时，应重新核对导数和取样标记。一直找不到满足阈值的点时，该时刻的代表密度可能缺失。}

\nextpage{shockspeed}{读取位置、速度与推算压强}
\subhead{先用位移复算速度}
本例的边缘从 18.5 µm 移到 30.5 µm，历时 12 ns。位移除以时间得到 1 µm/ns，即 1 km/s。打开\ui{冲击波速度}页，对照位置曲线的斜率和速度读数。要用这个方法计算物理速度，应选欧拉空间位置。材料坐标和网格编号的斜率有各自不同的含义。
\Figure{shock-velocity}{Ubuntu 上的速度页。位置随时间线性上升，底部速度轴设为 0.8 至 1.2 km/s，可直接读出整段轨迹保持在 1 km/s。}{82mm}
\Table{控件或输出}{怎样影响读数}{
速度页的 \ui{Frac}、\ui{Apply} & 轨迹至少有五个点时，显示的平滑曲线先用 LOWESS 拟合位置，在间隔 0.1 ns 的时间网上插值，再求导。记录出图时采用的 Frac。\\
速度时间范围 & 按 ns 填上下限，起点小于终点，区间内至少有两个轨迹点。清空两端可查看全程。\\
CSV 速度列 & 导出时先在已跟踪时刻数值求导。点数足够时，勾选常规参数中的“启用平滑”，程序会按指定窗口和多项式阶数，用 Savitzky--Golay 法对位置求导。它与速度页的 LOWESS 曲线可能给出不同读数。\\}
\subhead{怎样理解 Hugoniot 关系页}
查看此页时应选质量密度 \code{R}。程序使用跟踪速度、所选代表密度、初始压强 $P_0$（Pa）及初始密度 $\rho_0$（g/cm³）推算压强，并显示压强随压缩比的变化。该结果独立于文件中存储的总压强 \code{PT}。比较两次导出时应核对初始参数和密度取值方式。本练习的合成边缘可用来独立核算速度。

\nextpage{shockexport}{导出轨迹并逐列检查}
点击\ui{保存热图}可得到带轨迹的图像。\ui{保存轨迹}会为纳入时间区间的每个轨迹点写一行 CSV。练习时先清空速度时间上下限，关闭常规参数中的平滑，选最大密度，重新检测并保存为 \code{edge-max.csv}。改选跳变后密度，重新检测，另存为 \code{edge-jump.csv}。
\Table{CSV 字段}{练习数据的检查方法}{
\code{Time_Index}、\code{Time (ns)} & 共 13 行数据。索引从 0 开始，第 6 步的时间为 6.0000 ns。网格输出步号和物理时间各占一列。\\
\code{Position (um)}、\code{Position (cm)} & 第 6 步两列分别为 24.5000 µm 和 0.002450 cm。把厘米列乘以 10,000，应与微米列一致。\\
\code{Velocity (km/s)} & 本例为 1.0000。也可以用第 5 到第 7 步的位移除以 2 ns，再核一次。\\
\code{Max R (g/cm3)} 或 \code{Post-jump R (g/cm3)} & 列名随取值方式改变。第 6 步两个文件分别给出 1.500000 和 1.496654 g/cm³，位置都应为 24.5 µm。\\
压缩比和压强 & 压缩比为所取密度与 $\rho_0$ 之比。推算压强分别以 Pa 和 TPa 给出。密度缺失或压缩比小于等于 1 时，压强可能为空或显示为 \code{nan}。\\}
\subhead{轨迹中断、突然跳变或水平延伸时}
参考区间没有结果时，把范围移到肉眼可见的过渡区，并确认里面至少有三个网格中心。轨迹中断或斜率突变时，逐帧检查前后的剖面。找不到新位置时，如果上一帧的位置仍在有效范围内，跟踪器可能沿用该值。仅凭一段水平轨迹无法判断波前是否停止。2.3.1 在每个时刻最多输出一个位置，添加多个参考或勾选“允许多条冲击波”时同样如此。
\checknote{2.3.1 的摘要显示}{选择 \code{um} 时，结果摘要可能将起止位置都显示为 \code{0.0 um}。本例的位置可从轨迹叠加图和 CSV 的两列位置值核实。}
\subhead{复现上次分析}
窗口中的配置保存和加载按钮可恢复参考行、连续性约束、检测阈值及常规速度设置。参考行保存的是时间步索引，换文件后要对照实际时刻。另在 CSV 旁记录热图坐标与单位、代表密度模式、跳变后导数阈值及 LOWESS 的 Frac，下次打开时逐项核对。
